% Curva di riscaldamento di 1 kg di acqua, da ghiaccio a -20 °C a vapore a 120 °C.
% Calori (kJ): ghiaccio 2,1·20 = 42; fusione 334; acqua 4,2·100 = 420;
% ebollizione 2260; vapore 2,0·20 = 40. Cumulati: 42, 376, 796, 3056, 3096.
\begin{tikzpicture}
\begin{axis}[
  width=0.92\linewidth, height=5.2cm,
  axis lines=left, xmin=0, xmax=3300, ymin=-30, ymax=135,
  xlabel={calore fornito $Q$ (\si{\kilo\joule})}, ylabel={temperatura (\si{\celsius})},
  xtick={0,500,1000,1500,2000,2500,3000},
  ytick={-20,0,50,100},
  scaled x ticks=false,
  xticklabel style={/pgf/number format/1000 sep={}},
  tick label style={font=\small}, label style={font=\small},
  clip=false,
]
  \addplot[attenzione, very thick] coordinates {(0,-20) (42,0) (376,0) (796,100) (3056,100) (3096,120)};
  \draw[dashed, gray] (axis cs:0,0) -- (axis cs:42,0);
  \node[font=\scriptsize, anchor=south, align=center] at (axis cs:209,2) {fusione\\$Q = mL_f$};
  \node[font=\scriptsize, anchor=west] at (axis cs:600,45) {acqua: $Q = mc\,\Delta T$};
  \node[font=\scriptsize, anchor=south] at (axis cs:1926,102) {ebollizione: $Q = mL_v$};
  \node[font=\scriptsize, anchor=north west, align=left] at (axis cs:60,-8) {ghiaccio};
  \node[font=\scriptsize, anchor=east] at (axis cs:3070,120) {vapore};
\end{axis}
\end{tikzpicture}
